% --- Archon physics formalization source begin ---
% archon:physics
% archon:covers PhyXMiniProblems/problem_phyx_mini_0320.lean
% archon:source-report reports/phyx_mini/problem_phyx_mini_0320.source.json

\chapter{Physics problem phyx\_mini\_0320}
\label{ch:PhyXMiniProblems_problem_phyx_mini_0320}

\paragraph{Problem source.}
\#\# Physical scenario

Figure shows four tubes with lengths \$1.0\$ m or \$2.0\$ m, with one or two open ends as drawn. The third harmonic is set up in each tube, and some of the sound that escapes from them is detected by detector \$D\$, which moves directly away from the tubes.

\#\# Question

In terms of the speed of sound \$v\$, what speed must the detector have such that the detected frequency of the sound from tube 4 is equal to the tube's fundamental frequency?

\#\# Answer choices

- A: A: \$\textbackslash{}frac\{v\}\{2\}\$

- B: B: \$\textbackslash{}frac\{2v\}\{5\}\$

- C: C: \$\textbackslash{}frac\{v\}\{3\}\$

- D: D: \$\textbackslash{}frac\{2v\}\{3\}\$

\#\# Figure caption (auxiliary; use the image as primary evidence)

The image contains the following elements:

- Four horizontal bars on the left side, each with a gradient color that appears to be a light blue or blue-gray. The bars are numbered sequentially from top to bottom:

  1. The top bar is the shortest.
  2. The second bar is slightly longer.
  3. The third bar is longer than the previous two.
  4. The fourth bar is the longest.

- On the right side of the image, there is a circle filled with a gradient similar to the bars, transitioning from light to dark. This circle has a bold outline.

- An arrow extends from the circle, pointing to the right, indicating a direction or flow.

- The letter "D" is placed to the left of the circle and arrow, suggesting a label or designation for this part of the diagram.

The layout suggests a diagram possibly illustrating a process, hierarchy, or some variables measured against a scale, with the circle and arrow representing a directional or dynamic component.

\#\# Dataset metadata

Resonance and Harmonics; Multi-Formula Reasoning, Predictive Reasoning

\paragraph{Recorded answer/context.}
D: D: \$\textbackslash{}frac\{2v\}\{3\}\$

\paragraph{Figure/image path.}
/root/proposal\_for\_physic/science-mango/phyx\_mini\_run/phyx\_data/test\_image/320.png

\paragraph{Formalization target.}
create a compiling Lean file with sorry bodies at `PhyXMiniProblems/problem\_phyx\_mini\_0320.lean`. The Lean declarations must preserve the physical quantities, dimensions or dimensional roles, figure labels, governing-law hypotheses, and final relation expressed by this problem.
Use Mathlib/Physlib names found through LeanExplore where available. If a physics API is missing, introduce faithful local abstractions rather than scalar placeholder aliases.

\begin{theorem}[Physics formalization target]
\label{thm:physics:phyx_mini_0320:target}
\lean{PhyXMiniProblems.ProblemPhyXMini0320.tube4_fundamental_iff_detectorSpeed_twoThirds}
\uses{def:physics:phyx-mini-0320:phyxminiproblems-problemphyxmini0320-fourtubedopplersetup, def:physics:phyx-mini-0320:phyxminiproblems-problemphyxmini0320-matchesproblemandprimaryfigure, def:physics:phyx-mini-0320:phyxminiproblems-problemphyxmini0320-hasphysicalacousticparameters, def:physics:phyx-mini-0320:phyxminiproblems-problemphyxmini0320-satisfiestubemodeanddopplerlaws, def:physics:phyx-mini-0320:phyxminiproblems-problemphyxmini0320-detectorhearstube4fundamental, def:physics:phyx-mini-0320:phyxminiproblems-problemphyxmini0320-recordeddatasetanswer, def:physics:phyx-mini-0320:phyxminiproblems-problemphyxmini0320-detectorspeedagreeswithchoice}
The assigned autoformalize agent should translate this physics problem into Lean declarations in the covered file, with theorem and lemma proof bodies written as `by sorry`.
\end{theorem}
\begin{proof}
This is an autoformalization task, not a proof task. Produce statements that can later be proved by the physics prover without weakening the physical model.
\end{proof}

% --- Archon declaration topology begin ---
\section{Declaration topology}
\begin{definition}[Length Quantity]
  \label{def:physics:phyx-mini-0320:phyxminiproblems-problemphyxmini0320-lengthquantity}
  \lean{PhyXMiniProblems.ProblemPhyXMini0320.LengthQuantity}
  A nonnegative physical length, independent of the chosen unit system
\end{definition}
\begin{proof}
  This is the declared model definition of Length Quantity.
\end{proof}
\begin{definition}[Axial Position]
  \label{def:physics:phyx-mini-0320:phyxminiproblems-problemphyxmini0320-axialposition}
  \lean{PhyXMiniProblems.ProblemPhyXMini0320.AxialPosition}
  A signed axial position along the horizontal direction in the image
\end{definition}
\begin{proof}
  This is the declared model definition of Axial Position.
\end{proof}
\begin{definition}[Frequency Quantity]
  \label{def:physics:phyx-mini-0320:phyxminiproblems-problemphyxmini0320-frequencyquantity}
  \lean{PhyXMiniProblems.ProblemPhyXMini0320.FrequencyQuantity}
  A nonnegative physical acoustic frequency
\end{definition}
\begin{proof}
  This is the declared model definition of Frequency Quantity.
\end{proof}
\begin{definition}[Speed Quantity]
  \label{def:physics:phyx-mini-0320:phyxminiproblems-problemphyxmini0320-speedquantity}
  \lean{PhyXMiniProblems.ProblemPhyXMini0320.SpeedQuantity}
  A nonnegative physical speed with dimension length per time
\end{definition}
\begin{proof}
  This is the declared model definition of Speed Quantity.
\end{proof}
\begin{definition}[length Readout]
  \label{def:physics:phyx-mini-0320:phyxminiproblems-problemphyxmini0320-lengthreadout}
  \lean{PhyXMiniProblems.ProblemPhyXMini0320.lengthReadout}
  \uses{def:physics:phyx-mini-0320:phyxminiproblems-problemphyxmini0320-lengthquantity}
  Read a physical length in a selected length unit
\end{definition}
\begin{proof}
  This is the declared model definition of length Readout.
\end{proof}
\begin{definition}[length In Meters]
  \label{def:physics:phyx-mini-0320:phyxminiproblems-problemphyxmini0320-lengthinmeters}
  \lean{PhyXMiniProblems.ProblemPhyXMini0320.lengthInMeters}
  \uses{def:physics:phyx-mini-0320:phyxminiproblems-problemphyxmini0320-lengthquantity, def:physics:phyx-mini-0320:phyxminiproblems-problemphyxmini0320-lengthreadout}
  Metre readout of a physical length
\end{definition}
\begin{proof}
  This is the declared model definition of length In Meters.
\end{proof}
\begin{definition}[position Readout]
  \label{def:physics:phyx-mini-0320:phyxminiproblems-problemphyxmini0320-positionreadout}
  \lean{PhyXMiniProblems.ProblemPhyXMini0320.positionReadout}
  \uses{def:physics:phyx-mini-0320:phyxminiproblems-problemphyxmini0320-axialposition}
  Read a signed axial position in a selected length unit
\end{definition}
\begin{proof}
  This is the declared model definition of position Readout.
\end{proof}
\begin{definition}[position In Meters]
  \label{def:physics:phyx-mini-0320:phyxminiproblems-problemphyxmini0320-positioninmeters}
  \lean{PhyXMiniProblems.ProblemPhyXMini0320.positionInMeters}
  \uses{def:physics:phyx-mini-0320:phyxminiproblems-problemphyxmini0320-axialposition, def:physics:phyx-mini-0320:phyxminiproblems-problemphyxmini0320-positionreadout}
  Metre readout of a signed axial position
\end{definition}
\begin{proof}
  This is the declared model definition of position In Meters.
\end{proof}
\begin{definition}[frequency Readout]
  \label{def:physics:phyx-mini-0320:phyxminiproblems-problemphyxmini0320-frequencyreadout}
  \lean{PhyXMiniProblems.ProblemPhyXMini0320.frequencyReadout}
  \uses{def:physics:phyx-mini-0320:phyxminiproblems-problemphyxmini0320-frequencyquantity}
  Read a physical frequency in inverse units of a selected time unit
\end{definition}
\begin{proof}
  This is the declared model definition of frequency Readout.
\end{proof}
\begin{definition}[frequency In Hertz]
  \label{def:physics:phyx-mini-0320:phyxminiproblems-problemphyxmini0320-frequencyinhertz}
  \lean{PhyXMiniProblems.ProblemPhyXMini0320.frequencyInHertz}
  \uses{def:physics:phyx-mini-0320:phyxminiproblems-problemphyxmini0320-frequencyquantity, def:physics:phyx-mini-0320:phyxminiproblems-problemphyxmini0320-frequencyreadout}
  Hertz readout of a physical frequency
\end{definition}
\begin{proof}
  This is the declared model definition of frequency In Hertz.
\end{proof}
\begin{definition}[speed Readout]
  \label{def:physics:phyx-mini-0320:phyxminiproblems-problemphyxmini0320-speedreadout}
  \lean{PhyXMiniProblems.ProblemPhyXMini0320.speedReadout}
  \uses{def:physics:phyx-mini-0320:phyxminiproblems-problemphyxmini0320-speedquantity}
  Read a physical speed in selected length and time units
\end{definition}
\begin{proof}
  This is the declared model definition of speed Readout.
\end{proof}
\begin{definition}[speed In Meters Per Second]
  \label{def:physics:phyx-mini-0320:phyxminiproblems-problemphyxmini0320-speedinmeterspersecond}
  \lean{PhyXMiniProblems.ProblemPhyXMini0320.speedInMetersPerSecond}
  \uses{def:physics:phyx-mini-0320:phyxminiproblems-problemphyxmini0320-speedquantity, def:physics:phyx-mini-0320:phyxminiproblems-problemphyxmini0320-speedreadout}
  Metres-per-second readout of a physical speed
\end{definition}
\begin{proof}
  This is the declared model definition of speed In Meters Per Second.
\end{proof}
\begin{definition}[Tube Label]
  \label{def:physics:phyx-mini-0320:phyxminiproblems-problemphyxmini0320-tubelabel}
  \lean{PhyXMiniProblems.ProblemPhyXMini0320.TubeLabel}
  Tube numbers printed down the left side of the primary image
\end{definition}
\begin{proof}
  This is the declared model definition of Tube Label.
\end{proof}
\begin{definition}[Tube Endpoint]
  \label{def:physics:phyx-mini-0320:phyxminiproblems-problemphyxmini0320-tubeendpoint}
  \lean{PhyXMiniProblems.ProblemPhyXMini0320.TubeEndpoint}
  The two axial endpoints of any horizontal tube
\end{definition}
\begin{proof}
  This is the declared model definition of Tube Endpoint.
\end{proof}
\begin{definition}[Acoustic Boundary Condition]
  \label{def:physics:phyx-mini-0320:phyxminiproblems-problemphyxmini0320-acousticboundarycondition}
  \lean{PhyXMiniProblems.ProblemPhyXMini0320.AcousticBoundaryCondition}
  Acoustic boundary behavior indicated by the presence or absence of an end cap
\end{definition}
\begin{proof}
  This is the declared model definition of Acoustic Boundary Condition.
\end{proof}
\begin{definition}[Tube Motion]
  \label{def:physics:phyx-mini-0320:phyxminiproblems-problemphyxmini0320-tubemotion}
  \lean{PhyXMiniProblems.ProblemPhyXMini0320.TubeMotion}
  Motion state of a tube relative to the acoustic medium
\end{definition}
\begin{proof}
  This is the declared model definition of Tube Motion.
\end{proof}
\begin{definition}[Axial Direction]
  \label{def:physics:phyx-mini-0320:phyxminiproblems-problemphyxmini0320-axialdirection}
  \lean{PhyXMiniProblems.ProblemPhyXMini0320.AxialDirection}
  The two directions along the horizontal axis of the image
\end{definition}
\begin{proof}
  This is the declared model definition of Axial Direction.
\end{proof}
\begin{definition}[Detector Label]
  \label{def:physics:phyx-mini-0320:phyxminiproblems-problemphyxmini0320-detectorlabel}
  \lean{PhyXMiniProblems.ProblemPhyXMini0320.DetectorLabel}
  The detector label printed beneath the blue circular detector
\end{definition}
\begin{proof}
  This is the declared model definition of Detector Label.
\end{proof}
\begin{definition}[Acoustic Medium]
  \label{def:physics:phyx-mini-0320:phyxminiproblems-problemphyxmini0320-acousticmedium}
  \lean{PhyXMiniProblems.ProblemPhyXMini0320.AcousticMedium}
  The unspecified common gas in and around the four tubes
\end{definition}
\begin{proof}
  This is the declared model definition of Acoustic Medium.
\end{proof}
\begin{definition}[Acoustic Tube]
  \label{def:physics:phyx-mini-0320:phyxminiproblems-problemphyxmini0320-acoustictube}
  \lean{PhyXMiniProblems.ProblemPhyXMini0320.AcousticTube}
  \uses{def:physics:phyx-mini-0320:phyxminiproblems-problemphyxmini0320-lengthquantity, def:physics:phyx-mini-0320:phyxminiproblems-problemphyxmini0320-axialposition, def:physics:phyx-mini-0320:phyxminiproblems-problemphyxmini0320-tubeendpoint, def:physics:phyx-mini-0320:phyxminiproblems-problemphyxmini0320-acousticboundarycondition, def:physics:phyx-mini-0320:phyxminiproblems-problemphyxmini0320-tubemotion}
  Physical geometry and motion data associated with one pictured tube
\end{definition}
\begin{proof}
  This is the declared model definition of Acoustic Tube.
\end{proof}
\begin{definition}[Detector State]
  \label{def:physics:phyx-mini-0320:phyxminiproblems-problemphyxmini0320-detectorstate}
  \lean{PhyXMiniProblems.ProblemPhyXMini0320.DetectorState}
  \uses{def:physics:phyx-mini-0320:phyxminiproblems-problemphyxmini0320-axialposition, def:physics:phyx-mini-0320:phyxminiproblems-problemphyxmini0320-speedquantity, def:physics:phyx-mini-0320:phyxminiproblems-problemphyxmini0320-axialdirection, def:physics:phyx-mini-0320:phyxminiproblems-problemphyxmini0320-detectorlabel}
  Physical state of the moving detector D
\end{definition}
\begin{proof}
  This is the declared model definition of Detector State.
\end{proof}
\begin{definition}[Four Tube Doppler Setup]
  \label{def:physics:phyx-mini-0320:phyxminiproblems-problemphyxmini0320-fourtubedopplersetup}
  \lean{PhyXMiniProblems.ProblemPhyXMini0320.FourTubeDopplerSetup}
  \uses{def:physics:phyx-mini-0320:phyxminiproblems-problemphyxmini0320-lengthquantity, def:physics:phyx-mini-0320:phyxminiproblems-problemphyxmini0320-frequencyquantity, def:physics:phyx-mini-0320:phyxminiproblems-problemphyxmini0320-speedquantity, def:physics:phyx-mini-0320:phyxminiproblems-problemphyxmini0320-tubelabel, def:physics:phyx-mini-0320:phyxminiproblems-problemphyxmini0320-axialdirection, def:physics:phyx-mini-0320:phyxminiproblems-problemphyxmini0320-acousticmedium, def:physics:phyx-mini-0320:phyxminiproblems-problemphyxmini0320-acoustictube, def:physics:phyx-mini-0320:phyxminiproblems-problemphyxmini0320-detectorstate}
  The independent physical quantities in the four-tube experiment. The emitted and detected frequencies, fundamental frequencies, wavelengths, sound speed, and detector speed are independent fields. In particular, the setup does not assign the detector speed a multiple of the sound speed
\end{definition}
\begin{proof}
  This is the declared model definition of Four Tube Doppler Setup.
\end{proof}
\begin{definition}[Matches Problem And Primary Figure]
  \label{def:physics:phyx-mini-0320:phyxminiproblems-problemphyxmini0320-matchesproblemandprimaryfigure}
  \lean{PhyXMiniProblems.ProblemPhyXMini0320.MatchesProblemAndPrimaryFigure}
  \uses{def:physics:phyx-mini-0320:phyxminiproblems-problemphyxmini0320-lengthinmeters, def:physics:phyx-mini-0320:phyxminiproblems-problemphyxmini0320-positioninmeters, def:physics:phyx-mini-0320:phyxminiproblems-problemphyxmini0320-fourtubedopplersetup}
  Textual data and readouts from the primary image. The tube lengths and end conditions record all four depicted tubes even though only tube 4 enters the final calculation. Every right tube end lies to the left of detector D, the sound and detector move rightward, and the tubes are stationary; hence the detector moves directly away from each source. No detector-speed multiplier or detected-frequency tuning condition occurs in this structure
\end{definition}
\begin{proof}
  This is the declared model definition of Matches Problem And Primary Figure.
\end{proof}
\begin{definition}[Has Physical Acoustic Parameters]
  \label{def:physics:phyx-mini-0320:phyxminiproblems-problemphyxmini0320-hasphysicalacousticparameters}
  \lean{PhyXMiniProblems.ProblemPhyXMini0320.HasPhysicalAcousticParameters}
  \uses{def:physics:phyx-mini-0320:phyxminiproblems-problemphyxmini0320-lengthinmeters, def:physics:phyx-mini-0320:phyxminiproblems-problemphyxmini0320-frequencyinhertz, def:physics:phyx-mini-0320:phyxminiproblems-problemphyxmini0320-speedinmeterspersecond, def:physics:phyx-mini-0320:phyxminiproblems-problemphyxmini0320-fourtubedopplersetup}
  Positivity and subsonic conditions selecting the ordinary acoustic branch
\end{definition}
\begin{proof}
  This is the declared model definition of Has Physical Acoustic Parameters.
\end{proof}
\begin{definition}[Satisfies Tube Mode And Doppler Laws]
  \label{def:physics:phyx-mini-0320:phyxminiproblems-problemphyxmini0320-satisfiestubemodeanddopplerlaws}
  \lean{PhyXMiniProblems.ProblemPhyXMini0320.SatisfiesTubeModeAndDopplerLaws}
  \uses{def:physics:phyx-mini-0320:phyxminiproblems-problemphyxmini0320-lengthreadout, def:physics:phyx-mini-0320:phyxminiproblems-problemphyxmini0320-positionreadout, def:physics:phyx-mini-0320:phyxminiproblems-problemphyxmini0320-frequencyreadout, def:physics:phyx-mini-0320:phyxminiproblems-problemphyxmini0320-speedreadout, def:physics:phyx-mini-0320:phyxminiproblems-problemphyxmini0320-tubelabel, def:physics:phyx-mini-0320:phyxminiproblems-problemphyxmini0320-fourtubedopplersetup}
  Governing relations for the standing modes and the sound heard by D. An open--open mode obeys 2 L = n λ. A mode with exactly one closed end obeys 4 L = n λ, with odd n. A harmonic's emitted frequency is n times its fundamental frequency. Sound propagation obeys v = f λ. For a stationary source and a subsonic observer receding in the direction of propagation, the detected frequency is f (v-u)/v. The laws are stated in arbitrary compatible units.
\end{definition}
\begin{proof}
  This is the declared model definition of Satisfies Tube Mode And Doppler Laws.
\end{proof}
\begin{definition}[Detector Hears Tube4 Fundamental]
  \label{def:physics:phyx-mini-0320:phyxminiproblems-problemphyxmini0320-detectorhearstube4fundamental}
  \lean{PhyXMiniProblems.ProblemPhyXMini0320.DetectorHearsTube4Fundamental}
  \uses{def:physics:phyx-mini-0320:phyxminiproblems-problemphyxmini0320-fourtubedopplersetup}
  The condition in the question: the sound from tube 4 is detected at tube 4's fundamental frequency. This is a proposition about two independent physical frequency fields, not a definition of either frequency or of detector speed
\end{definition}
\begin{proof}
  This is the declared model definition of Detector Hears Tube4 Fundamental.
\end{proof}
\begin{definition}[Answer Choice]
  \label{def:physics:phyx-mini-0320:phyxminiproblems-problemphyxmini0320-answerchoice}
  \lean{PhyXMiniProblems.ProblemPhyXMini0320.AnswerChoice}
  Labels of the four displayed speed choices
\end{definition}
\begin{proof}
  This is the declared model definition of Answer Choice.
\end{proof}
\begin{definition}[sound Speed Multiplier]
  \label{def:physics:phyx-mini-0320:phyxminiproblems-problemphyxmini0320-answerchoice-soundspeedmultiplier}
  \lean{PhyXMiniProblems.ProblemPhyXMini0320.AnswerChoice.soundSpeedMultiplier}
  \uses{def:physics:phyx-mini-0320:phyxminiproblems-problemphyxmini0320-answerchoice}
  Dimensionless multiplier of the sound speed printed beside each choice
\end{definition}
\begin{proof}
  This is the declared model definition of sound Speed Multiplier.
\end{proof}
\begin{definition}[recorded Dataset Answer]
  \label{def:physics:phyx-mini-0320:phyxminiproblems-problemphyxmini0320-recordeddatasetanswer}
  \lean{PhyXMiniProblems.ProblemPhyXMini0320.recordedDatasetAnswer}
  \uses{def:physics:phyx-mini-0320:phyxminiproblems-problemphyxmini0320-answerchoice}
  The answer label recorded in the supplied dataset metadata
\end{definition}
\begin{proof}
  This is the declared model definition of recorded Dataset Answer.
\end{proof}
\begin{definition}[Detector Speed Agrees With Choice]
  \label{def:physics:phyx-mini-0320:phyxminiproblems-problemphyxmini0320-detectorspeedagreeswithchoice}
  \lean{PhyXMiniProblems.ProblemPhyXMini0320.DetectorSpeedAgreesWithChoice}
  \uses{def:physics:phyx-mini-0320:phyxminiproblems-problemphyxmini0320-fourtubedopplersetup, def:physics:phyx-mini-0320:phyxminiproblems-problemphyxmini0320-answerchoice}
  A displayed choice agrees with the detector's physical speed
\end{definition}
\begin{proof}
  This is the declared model definition of Detector Speed Agrees With Choice.
\end{proof}
% --- Archon declaration topology end ---
% --- Archon physics formalization source end ---
